\begin{figure}[tb]
\centering
\includegraphics[clip, trim=0cm 0.2cm 0cm 0cm, width=\textwidth]{figures/usage.pdf}
\caption{Illustration of the applications of Step-Audio 2 across various speech conversation scenarios.}
\label{fig:usage}
\end{figure}
\section{Related Work}

\subsection{Speech and audio understanding}

Recent advances in large language models (LLMs)~\cite{bai2023qwentechnicalreport,grattafiori2024llama3herdmodels,openai2023gpt4,openai2022chatgpt} have extended their application to a wide range of speech and audio understanding tasks, such as audio captioning, sound event detection, automatic speech recognition, audio classification, and audio-driven creative generation.
A prevalent approach~\cite{chu2024qwen2,deshmukh2023pengi,gong2023joint,gong2023listen,ParalinGPT,tang2023salmonn} involves pairing speech encoders with lightweight, trainable adapters that project audio features into a textual embedding space compatible with LLMs.
Building on this foundation, recent studies have further explored how to incorporate paralinguistic information such as emotion, intonation and speaker style, enabling LLMs to move beyond pure linguistic comprehension. For instance, ParalinGPT~\cite{ParalinGPT} focuses on enhancing a powerful text-based language model by integrating continuous speech embeddings, enabling it to capture paralinguistic signals such as emotion and prosody. SALMONN~\cite{tang2023salmonn} adopts a multi-modal strategy by freezing speech encoders Whisper~\cite{radford2023robust} and BEATs~\cite{BEATs}, and connecting their outputs to an LLM via a window-level Q-Former, enabling joint modeling of linguistic and acoustic features. Seed-ASR~\cite{Seedasr} integrates LUISE-based speech representations with instructions and context, using context-aware SFT to capture semantic information. AudioPaLM~\cite{rubenstein2023audiopalm} combines PaLM-2~\cite{anil2023palm2technicalreport} and AudioLM~\cite{borsos2023audiolm}, unifying linguistic knowledge with paralinguistic features like speaker identity and intonation. LLM-based approaches~\cite{chu2023qwen,chu2024qwen2} increasingly rely on pretrained audio encoders such as Wav2Vec~\cite{wav2vec2}, HuBERT~\cite{HuBERT}, Whisper~\cite{radford2023robust}, and WavLM~\cite{WavLM} to extract rich semantic representations from speech. At the same time, the extensive text knowledge and contextual reasoning capabilities stored in LLMs can provide valuable semantic guidance for understanding tasks. 




\subsection{Text-to-speech synthesis}

Text-to-Speech (TTS) technology has made remarkable strides in recent years, evolving from traditional concatenative and statistical parametric approaches~\cite{deepvoice3,fastspeech2,tacotron2,wang2017tacotron} to codec-based TTS systems.
Codec language models leverage a speech codec to extract discrete representations of speech~\cite{encodc,defossez2024moshi,ji2024wavtokenizer,snac,,xin2024bigcodec,soundstream,zhang2025distinctive,zhang2023speechtokenizer} and utilize either autoregressive~\cite{ding2025kimi,zhang2023speechgpt} or masked language models~\cite{wang2024maskgct} to predict the corresponding speech tokens. These tokens are then synthesized into waveforms using codec vocoders. VALL-E~\cite{wang2023neural} marked a significant breakthrough in this area. It uses an autoregressive model to generate coarse codec codes, followed by a non-autoregressive model for the fine codes. Unlike VALL-E, which predicts acoustic tokens from phonemes and requires transcripts, SPEAR-TTS~\cite{SPEARTTS} uses a two-stage architecture with self-supervised audio prompts to clone unseen voices from just 3 seconds of speech. SparkTTS~\cite{wang2025spark} introduces BiCodec, a single-stream speech codec that encodes linguistic content as compact semantic tokens and speaker characteristics as fixed-length global tokens. Instead of relying on non-autoregressive models to predict residual discrete codes, methods like TorToiseTTS~\cite{TorToiseTTS}, CosyVoice~\cite{du2024CosyVoice}, CosyVoice~2~\cite{du2024cosyvoice2}, MiniMax-Speech~\cite{zhang2025minimaxspeech} and SEED-TTS~\cite{seedtts} adopt diffusion or flow-matching techniques as a second stage to reconstruct mel-spectrograms or continuous representations enriched with fine-grained acoustic and semantic details. Recent work, Kimi-Audio~\cite{ding2025kimi} combines Whisper features and semantic tokens for efficient modeling, with dual heads and a flow-matching detokenizer plus BigVGAN~\cite{lee2023bigvgan} for low-latency, expressive synthesis.


\subsection{Speech-to-speech translation}
Speech-to-speech translation (S2ST) is a crucial technology for eliminating communication barriers across languages. Traditional S2ST systems~\cite{wahlster2013verbmobil,wu2016google} typically adopt a cascaded pipeline consisting of automatic speech recognition (ASR), machine translation (MT), and TTS modules. Earlier studies~\cite{jia2019direct,jia2022translatotron,le2024transvip,lee2021textless} have shifted toward direct approaches that bypass intermediate textual representations, aiming for lower latency and better preservation of prosody and speaker characteristics. Two main types of direct S2ST methods have emerged, which are known as speech-to-spectrogram translation and speech-to-unit translation. Both directly generate target speech representations from the source speech without relying on textual transcriptions. A representative of the former is Translatotron~\cite{jia2019direct}, the first end-to-end model to translate source speech directly into target spectrograms. Translatotron 2~\cite{jia2022translatotron}, further improves translation quality through a two-pass decoding mechanism~\cite{lee2021textless}. In contrast, speech-to-unit models predict discrete acoustic tokens rather than spectrograms, which are typically extracted using self-supervised speech encoders such as HuBERT~\cite{HuBERT} or WavLM~\cite{WavLM}. For instance, TransVIP~\cite{le2024transvip} employs a joint encoder-decoder architecture that first generates target text and residual vector quantization (RVQ) codes in the initial layer, followed by a non-causal language model that refines RVQ predictions in subsequent layers.


\subsection{Speech-to-text and speech-to-speech conversation}

Based on whether the LLM can directly understand and generate speech representations, existing systems can be categorized into end-to-end large audio language models and cascaded large audio language models. The former directly models audio inputs and outputs within a unified framework, while the latter relies on a modular pipeline involving separate ASR, LLM, and TTS components. Traditional speech-to-text and speech-to-speech systems typically adopt a cascaded architecture, as exemplified by AudioGPT~\cite{huang2024audiogpt} and Spoken-LLM~\cite{lin2024advancing}. However, the ASR + LLM + TTS pipeline incurs high latency and modular mismatches. This has spurred interest in unified end-to-end architectures for faster and more seamless integration. A major milestone in this direction is GPT-4o~\cite{hurst2024gpt}, which supports direct end-to-end speech interaction without requiring intermediate textual conversions. Recently, several new end-to-end LALMs~\cite{defossez2024moshi,fang2024llama,xie2024mini,xie2024mini2,xu2025qwen25omnitechnicalreport} for speech-to-speech conversation have emerged. For instance, Moshi~\cite{defossez2024moshi} improves efficiency with an RQ-Transformer that generates text and audio tokens simultaneously. Similarly, Mini-Omni~\cite{xie2024mini} generates speech and text responses in parallel, following a strategy similar to MusicGen~\cite{musicgen}, which enables lower first-token latency compared to interleaved generation designs. LUCY~\cite{gao2025lucylinguisticunderstandingcontrol} builds on the Mini-Omni architecture with enhancements for emotional expressiveness, naturalness, and informativeness in speech generation. It utilizes curated synthetic data and optimizes the training and decoding pipelines to handle multi-turn dialogue and function-call scenarios. Mini-Omni2~\cite{xie2024mini} further extends Mini-Omni framework by integrating multimodal understanding and full-duplex interaction capabilities. LLaMA-Omni~\cite{fang2024llama} introduces a streaming, non-autoregressive speech decoder based on Connectionist Temporal Classification, enabling direct and efficient generation of discrete audio tokens without relying on step-by-step prediction. Freeze-Omni~\cite{wang2024freeze}, on the other hand, freezes the LLM parameters during training, preserving its original capabilities while achieving low-latency speech-to-speech interaction through streaming and decoder integration. Qwen2.5-Omni~\cite{xu2025qwen25omnitechnicalreport} supports multimodal input and simultaneous text-speech output via a thinker-talker architecture, using TMRoPE for improved audio-visual synchronization through explicit temporal encoding.








